\documentclass[14pt]{beamer}
\usetheme{Berkeley}
\usepackage[utf8]{inputenc}
\usepackage[ngerman]{babel}
\usepackage[T1]{fontenc}
\usepackage{amsmath}
\usepackage{amsfonts}
\usepackage{amssymb}
\usepackage{graphicx}
%\usepackage{svg}
%\usepackage{pdfplots}
%\usepackage{pdfplotstable}
\graphicspath{{images/}}
\author{AUTHOR}
\title{TITLE}
\usepackage{mathptmx}
%\setbeamercovered{transparent} 
%\setbeamertemplate{navigation symbols}{} 
\logo{\includegraphics[scale=0.4]{logo.png}} 
\institute{Herder Gymnasium Berlin} 
\date{} 
%\subject{test} 
\begin{document}

\begin{frame}
\titlepage
\end{frame}

\begin{frame}
\tableofcontents
\end{frame}

\section{Was ist eine Dichtefunktion?}
\begin{frame}{Wozu eine Dichtefunktion?}
\textbf{Bisherige Berechnung}:

{\small $$ P(X \in\mathbb{A}) \, = \, \sum_{x\in A} \rho(x) $$}
\begin{itemize}

\item eignet sich nur für diskrete Wahrscheinlichkeitsverteilungen:

\begin{itemize}
\only<1>{
\item Binomial
\item Hypergeometrisch
\item Poisson
}
\only<2>{
\item[$\Rightarrow$] keine Aussage möglich für stetige Verteilungen, da $P(X \, = \, x) = 0 \, \forall X \in \mathbb{A}$
}
\end{itemize}
\end{itemize}

\end{frame}

\begin{frame}{Definition}
\textbf{Wahrscheinlichkeit als Flächeninhalt:}
$\varphi \colon \mathbb{R} \rightarrow [0,\infty)$ mit
$$P(X \in \mathbb{A})=\int_{\mathbb{A}} f(x)\,\mathrm dx$$
ist die \textbf{Wahrscheinlichkeitsdichte} der Verteilung von $X$.
\end{frame}

\begin{frame}{Eigenschaften}
Eigenschaften:
\begin{itemize}
\item $\forall x \in \mathbb{R}:\, f(x)\geq 0$
\item $\int_{-\infty}^\infty f(x)\,{\mathrm d}x = 1$
\end{itemize}
Eine Funktion $f$ ist genau dann eine Dichtefunktion, wenn diese Eigenschaften erfüllt sind.
\end{frame}

\begin{frame}{Berechnen mithilfe der Dichtefunktion}
Wahrscheinlichkeit für ein Intervall $[a,b]$:
$$P(X\in [a,b]) = \int_a^b f(x)\, {\rm d}x = F(b)-F(a).$$
$F$ ist die sogenannte \textbf{kumultative Verteilungsfunktion} mit:
$$F(x) \, = \, P(X \le x)$$
\end{frame}

\begin{frame}{Suche der Funktion}
\textbf{Ziel: Annäherung an die Binomialverteilung:}
\pause
\begin{itemize}
\item Verschieben des Erwartungswertes, sodass $\mu=0$:\\
$k \to k-\mu$
\item Vereinheitlichen der Breite:\\
$k-\mu \to \frac{k-\mu}{\sigma}$
\pause
\item \glqq Wiederherstellen\grqq\ des Flächeninhalts:\\
$B(n;p;k) \to \sigma \cdot B(n;p;k)$
\end{itemize}

\end{frame}
\begin{frame}{Zwischenstand}
$$
\sigma \cdot B(n;p;k) =: \varphi\ \left(\ \frac{k-\mu}{\sigma}\right)
$$
$$
\Leftrightarrow
\begin{mathbf}
B(n;p;k) =  \frac{1}{\sigma} \cdot  \varphi\ \left(\ \frac{k-\mu}{\sigma}\right)
\end{mathbf}
$$
\end{frame}






\section{Was sagt der Grenzwertsatz aus?} % (Bedeutung)
\subsection{Beweis} %  od. Plausibilitätsprüfung
\section{Gauß} %Hinführung auf Gauß'sche Formel (e^(-1/2x^2)), so wie es Herr Steinkrauß erläutert hat (oder auch anders, aber in diesem Detailgrad)
\section{Aufgaben zur lokalen Näherung}
\end{document}



Was ich dann (nachfolgend) mache:
- Was sagt das Integral aus?
- Wie bilde ich überhaupt das Integral?
- Wie gucke ich die Werte nach? (Tabellen)
+ ein paar Aufgaben 